\documentclass[UTF8,12pt]{ctexart}
\usepackage[paperwidth=240mm,paperheight=200mm,margin=14mm,top=8mm]{geometry}
\usepackage{amsmath,amssymb,mathtools}
\usepackage{xcolor}
\pagestyle{empty}
\setlength{\parindent}{0pt}
\setlength{\parskip}{0pt}
\newcommand{\qn}[1]{\textbf{\large #1}\ }
\xeCJKDeclareCharClass{CJK}{"2460 -> "24FF}
\begin{document}
\qn{1993.1}
$\lim\limits_{x\to 0^+}x\ln x=$______.

\vfill
\newpage
\qn{1993.2}
函数 $y=y(x)$ 由方程 $\sin(x^2+y^2)+\mathrm{e}^x-xy^2=0$ 所确定，则 $\dfrac{\mathrm{d}y}{\mathrm{d}x}=$______.

\vfill
\newpage
\qn{1993.3}
设 $F(x)=\int_1^x\left(2-\dfrac{1}{\sqrt{t}}\right)\mathrm{d}t\ (x>0)$，则函数 $F(x)$ 的单调减少区间是______.

\vfill
\newpage
\qn{1993.4}
$\int\dfrac{\tan x}{\sqrt{\cos x}}\,\mathrm{d}x=$______.

\vfill
\newpage
\qn{1993.5}
已知曲线 $y=f(x)$ 过点 $\left(0,-\dfrac{1}{2}\right)$，且其上任一点 $(x,y)$ 处的切线斜率为 $x\ln(1+x^2)$，则 $f(x)=$______.

\vfill
\newpage
\qn{1993.6}
当 $x\to 0$ 时，变量 $\dfrac{1}{x^2}\sin\dfrac{1}{x}$ 是（　　）
\\
\noindent (A) 无穷小.　　(B) 无穷大.
\\
\noindent (C) 有界的，但不是无穷小.　　(D) 无界的，但不是无穷大.

\vfill
\newpage
\qn{1993.7}
设 $f(x)=\begin{cases}\dfrac{|x^2-1|}{x-1}, & x\neq 1,\\ 2, & x=1,\end{cases}$ 则在点 $x=1$ 处函数 $f(x)$（　　）
\\
\noindent (A) 不连续.　　(B) 连续，但不可导.
\\
\noindent (C) 可导，但导数不连续.　　(D) 可导，且导数连续.

\vfill
\newpage
\qn{1993.8}
已知 $f(x)=\begin{cases}x^2, & 0\leqslant x<1,\\ 1, & 1\leqslant x\leqslant 2,\end{cases}$ 设 $F(x)=\int_1^x f(t)\,\mathrm{d}t\ (0\leqslant x\leqslant 2)$，则 $F(x)$ 为（　　）
\\
\noindent (A) $\begin{cases}\dfrac{1}{3}x^3, & 0\leqslant x<1,\\ x, & 1\leqslant x\leqslant 2.\end{cases}$
\\
\noindent (B) $\begin{cases}\dfrac{1}{3}x^3-\dfrac{1}{3}, & 0\leqslant x<1,\\ x, & 1\leqslant x\leqslant 2.\end{cases}$
\\
\noindent (C) $\begin{cases}\dfrac{1}{3}x^3, & 0\leqslant x<1,\\ x-1, & 1\leqslant x\leqslant 2.\end{cases}$
\\
\noindent (D) $\begin{cases}\dfrac{1}{3}x^3-\dfrac{1}{3}, & 0\leqslant x<1,\\ x-1, & 1\leqslant x\leqslant 2.\end{cases}$

\vfill
\newpage
\qn{1993.9}
设常数 $k>0$，函数 $f(x)=\ln x-\dfrac{x}{\mathrm{e}}+k$ 在 $(0,+\infty)$ 内的零点个数为（　　）
\\
\noindent (A) 3.　　(B) 2.
\\
\noindent (C) 1.　　(D) 0.

\vfill
\newpage
\qn{1993.10}
若 $f(x)=-f(-x)$，在 $(0,+\infty)$ 内 $f'(x)>0$，$f''(x)>0$，则 $f(x)$ 在 $(-\infty,0)$ 内（　　）
\\
\noindent (A) $f'(x)<0$，$f''(x)<0$.　　(B) $f'(x)<0$，$f''(x)>0$.
\\
\noindent (C) $f'(x)>0$，$f''(x)<0$.　　(D) $f'(x)>0$，$f''(x)>0$.

\vfill
\newpage
\qn{1993.11}
设 $y=\sin[f(x^2)]$，其中 $f$ 具有二阶导数，求 $\dfrac{\mathrm{d}^2y}{\mathrm{d}x^2}$.

\vfill
\newpage
\qn{1993.12}
求 $\lim\limits_{x\to-\infty}x(\sqrt{x^2+100}+x)$.

\vfill
\newpage
\qn{1993.13}
求 $\int_0^{\frac{\pi}{4}}\dfrac{x}{1+\cos 2x}\,\mathrm{d}x$.

\vfill
\newpage
\qn{1993.14}
求 $\int_0^{+\infty}\dfrac{x}{(1+x)^3}\,\mathrm{d}x$.

\vfill
\newpage
\qn{1993.15}
求微分方程 $(x^2-1)\,\mathrm{d}y+(2xy-\cos x)\,\mathrm{d}x=0$ 满足初值条件 $y(0)=1$ 的特解.
设二阶常系数线性微分方程 $y''+\alpha y'+\beta y=\gamma\mathrm{e}^x$ 的一个特解为 $y=\mathrm{e}^{2x}+(1+x)\mathrm{e}^x$，试确定常数 $\alpha,\beta,\gamma$，并求该方程的通解.
设平面图形 $A$ 由 $x^2+y^2\leqslant 2x$ 与 $y\geqslant x$ 所确定，求图形 $A$ 绕直线 $x=2$ 旋转一周所得旋转体的体积.
作半径为 $r$ 的球的外切正圆锥，问此圆锥的高 $h$ 为何值时，其体积 $V$ 最小，并求出该最小值.
设 $x>0$，常数 $a>\mathrm{e}$. 证明：$(a+x)^a<a^{a+x}$.
设 $f'(x)$ 在 $[0,a]$ 上连续，且 $f(0)=0$，证明：$\left|\int_0^a f(x)\,\mathrm{d}x\right|\leqslant\dfrac{Ma^2}{2}$，其中 $M=\max\limits_{0\leqslant x\leqslant a}|f'(x)|$.

\vfill
\newpage
\end{document}